\section*{Resumo}

Este trabalho apresenta uma avaliação comparativa de técnicas de extração de características para verificação biométrica por impressões digitais, contrastando as arquiteturas de aprendizado profundo \textit{DeepPrint} (vetor unidimensional de tamanho fixo) e \textit{FLARE} (descritor denso tridimensional com alinhamentos por regressão e votação) com o método clássico \textit{Minutia Cylinder-Code} (MCC). Os experimentos foram conduzidos sobre os subconjuntos Parte A das competições FVC2000, FVC2002 e FVC2004, além das bases ópticas CrossMatch e U.are.U da Neurotechnology, avaliando-se as taxas FMR, FNMR e o EER no ponto de mínima diferença. Os resultados indicam a superioridade expressiva do FLARE, que obteve EER médio geral de 0,64\% na abordagem por votação, beneficiado pela preservação de relações espaciais e supressão explícita de fundo. O DeepPrint alcançou EER médio de 5,92\%, apresentando maior sensibilidade a ruídos severos. O MCC obteve desempenho competitivo em bases ópticas de alta qualidade (EER médio de 2,02\% na Neurotechnology e 2,69\% no FVC2000 DB1-A), superando o DeepPrint nesses cenários, porém com degradação em amostras ruidosas. Por fim, constatou-se que o emprego de etapas prévias de realce (UNetEnh e PriorEnh) não resultou em aprimoramentos consistentes de desempenho nas arquiteturas avaliadas.